\documentclass[utf8]{FrontiersinHarvard}
\usepackage{url,hyperref,lineno,microtype,subcaption}
\usepackage[onehalfspacing]{setspace}
\usepackage{amsmath,amssymb,booktabs,tabularx,array,graphicx}
\linenumbers
\providecommand{\tightlist}{}

\def\keyFont{\fontsize{8}{11}\helveticabold}
\def\firstAuthorLast{Speiser}
\def\Authors{Leonard Speiser\,$^{1,*}$}
\def\Address{$^{1}$Horizon 3, Independent Research, Los Altos, CA, United States}
\def\corrAuthor{Leonard Speiser}
\def\corrEmail{leonard@horizon3.net}

\begin{document}
\onecolumn
\firstpage{1}
\title[$\Sigma$-Gravity]{$\Sigma$-Gravity: Coherence-Motivated Gravitational Enhancement in Galaxies and Clusters}
\author[\firstAuthorLast]{\Authors}
\address{}
\correspondance{}
\extraAuth{}
\maketitle

\begin{center}
\small Manuscript ID 1866133 \quad Word count: 3210 \quad Figures: 4 \quad Tables: 2
\end{center}

\begin{abstract}
\section{}
The dynamics of galaxies and galaxy clusters exceed the predictions obtained from directly observed baryonic matter, a discrepancy conventionally modeled with nonbaryonic dark matter or, at galaxy scales, phenomenological modifications such as modified Newtonian dynamics (MOND). We investigate \(\Sigma\)-Gravity as an alternative nonrelativistic phenomenological parameterization of selected missing-mass observations. Its empirical response is \(\Sigma(g_N,B)=1+B h(g_N)\), where \(h(g_N)=\sqrt{g^\dagger/g_N}\,g^\dagger/(g^\dagger+g_N)\) and \(g^\dagger=cH_0/(4\sqrt{\pi})\); the locked galaxy implementation supplements this acceleration dependence with a bounded factor motivated by rotational support. For a disk-dominated sample of 164 SPARC galaxies, using fixed stellar mass-to-light ratios and no per-galaxy halo fitting, the mean velocity root-mean-square residual is \(16.366\ {\rm km\,s^{-1}}\) for \(\Sigma\)-Gravity and \(16.056\ {\rm km\,s^{-1}}\) for the MOND prescription tested with the same baryonic inputs. Their paired difference is statistically consistent with comparable aggregate performance. Relative to the acceleration-only form, the bounded factor improves the mean RMS by \(0.517\ {\rm km\,s^{-1}}\), with a galaxy-bootstrap 95\% interval of \([0.242,0.795]\ {\rm km\,s^{-1}}\). Thus the factor makes a measurable contribution within the locked predictor, although its endogenous construction does not establish independently measured coherence. At cluster scales, an amplitude calibrated on 42 Fox et al.~strong-lensing systems gives a median predicted-to-observed aperture-mass ratio of 0.987 under a simplified baryon prescription. A no-refit check using 73 radial measurements in 17 non-overlapping CLASH clusters gives a median ratio of 1.32 and reveals a radius-dependent bias. An action-based QUMOND embedding exists for independently specified \(B\). Together, these results identify a low-parameter galaxy response competitive with MOND and expose a falsifiable limitation of its fixed cluster amplitude. The physical origin of \(B\), including possible kinematic coherence or another property of baryonic organization, remains open; no relativistic or cosmological extension is developed here.

\tiny
\keyFont{\section{Keywords:} galaxy rotation curves, radial acceleration relation, modified gravity, dark matter, galaxy clusters, gravitational lensing, SPARC, phenomenology}
\end{abstract}

\section{Introduction}\label{introduction}



Galaxy rotation curves, galaxy-cluster dynamics, gravitational lensing, the cosmic microwave background, and large-scale structure collectively motivate the standard \(\Lambda\)CDM framework \citep{planck2020}. Within that framework, nonbaryonic cold dark matter supplies the dominant gravitating mass in galaxies and clusters and plays a central role in cosmological structure formation. The microscopic identity of dark matter nevertheless remains unknown, and individual galaxy analyses commonly require system-specific halo parameters. These facts do not negate the cosmological evidence for dark matter, but they motivate tests of compact empirical relationships between baryonic structure and the observed gravitational discrepancy.

MOND provides the best-known alternative phenomenology at galaxy scales \citep{milgrom1983,sanders2002,famaey2012}. It relates the effective and Newtonian accelerations through a universal acceleration scale and organizes many disk-galaxy observations, including the radial acceleration relation \citep{mcgaugh2016}. Residual mass discrepancies in galaxy clusters and the need for additional structure in relativistic and cosmological realizations motivate the investigation of other phenomenological responses.

\(\Sigma\)-Gravity is evaluated here as an alternative phenomenological description of selected missing-mass observations. Its galaxy implementation combines a fixed acceleration response with a bounded factor motivated by rotational support and uses no per-galaxy response amplitude. For a disk rotation curve, this prescription can be compared directly with a MOND interpolation function or a per-galaxy dark-matter halo fit. The term ``alternative'' refers to this empirical galaxy-scale choice; relativistic and cosmological extensions lie outside the present scope.

The observations analyzed here identify the combined response amplitude \(B\), rather than independently separating an amplitude and a physical coherence variable. The implemented galaxy factor is calculated from the model-predicted velocity rather than independent phase-space measurements, and assigning every cluster the same nominal path length identifies one cluster amplitude rather than a length exponent. Coherence, source geometry, matter currents, or another property of baryonic organization therefore remain physical motivations to be tested rather than established ingredients of the empirical law.

This study defines the empirical response and its identifiable parameters, measures the contribution of the bounded galaxy factor through an acceleration-only ablation, provides a nonrelativistic action embedding, compares the locked galaxy predictor with MOND under common baryonic assumptions, and distinguishes cluster calibration from no-refit evaluation. These analyses preserve a falsifiable phenomenology while specifying what a deeper theory would need to derive.

\section{Materials and methods}\label{materials-and-methods}

\subsection{Canonical empirical response and identifiability}\label{canonical-empirical-response-and-identifiability}

Let \(\Phi_N\) be the Newtonian potential of the baryonic density \(\rho_b\),

\[
\nabla^2\Phi_N=4\pi G\rho_b,
\tag{1}
\]

and define \(g_N=|\nabla\Phi_N|\). The canonical empirical enhancement is

\[
\Sigma(g_N,B)=1+B\,h(g_N),
\tag{2}
\]

where

\[
h(g_N)=\sqrt{\frac{g^\dagger}{g_N}}
\frac{g^\dagger}{g^\dagger+g_N},
\qquad
g^\dagger=\frac{cH_0}{4\sqrt{\pi}}
\simeq9.60\times10^{-11}\ {\rm m\,s^{-2}} .
\tag{3}
\]

For \(g_N\ll g^\dagger\), \(h\sim\sqrt{g^\dagger/g_N}\); for \(g_N\gg g^\dagger\), \(h\rightarrow0\). In the deep-acceleration point-mass limit,

\[
V^4\rightarrow B^2GM_b g^\dagger .
\tag{4}
\]

The baryonic Tully--Fisher normalization therefore constrains \(B^2g^\dagger\), not \(B\) and \(g^\dagger\) independently. The observations considered here identify only the combined amplitude \(B\); separate interpretations in terms of \(B=A\mathcal C\) are therefore not assumed, and possible meanings of \(A\) and \(\mathcal C\) are deferred to the Discussion.

\subsection{Galaxy implementation}\label{galaxy-implementation}

The locked galaxy implementation uses the bounded fixed-point factor

\[
F(V)=\frac{V^2}{V^2+\sigma^2},
\qquad \sigma=20\ {\rm km\,s^{-1}},
\tag{5}
\]

with

\[
B_{\rm gal}=A_0F(V_\Sigma),\qquad
V_\Sigma=V_{\rm bar}\sqrt{1+B_{\rm gal}h(g_N)},\qquad
A_0=e^{1/(2\pi)} .
\tag{6}
\]

The iteration uses \(V_\Sigma\), not the observed velocity, and therefore does not directly fit the observed outcome. However, \(F(V_\Sigma)\) is endogenous to the prediction rather than an independently measured phase-space quantity. It is treated here as a phenomenological regularizer, not as a measured covariant coherence scalar.

The catalog photometric scale length \(R_d\), a radial window \(r/(R_d/2\pi+r)\), the reference length \(L_0\), and the exponent \(n\) do not enter this locked galaxy predictor and are not used to obtain the SPARC results reported below. The results therefore test the fixed-point implementation in Equation (6); incorporating a scale-length-dependent window would define a distinct extension requiring independent specification and evaluation.

A second velocity moment can schematically be decomposed in a stationary axisymmetric system as \(\langle v^2\rangle\simeq v_{\rm rot}^2+\sigma^2\). The ratio \(v_{\rm rot}^2/\langle v^2\rangle\) motivates a bounded measure of rotational support, but it is not a derivation of Equation (5). The decomposition is frame dependent; anisotropic systems require the dispersion tensor; and the implemented velocity is predicted rather than measured. No unique covariant reduction is claimed.

\subsection{Nonrelativistic action embedding}\label{nonrelativistic-action-embedding}

An action-based embedding can be written in the QUMOND framework \citep{milgrom2010} when \(B\) is independently specified. Define

\[
z=\frac{|\nabla\Phi_N|^2}{(g^\dagger)^2}
\tag{7}
\]

and

\[
Q(z,B)=z+4B\left[z^{1/4}-\arctan\!\left(z^{1/4}\right)\right].
\tag{8}
\]

Then

\[
Q_z=1+B\frac{z^{-1/4}}{1+\sqrt z}
=1+B h(g_N).
\tag{9}
\]

The gravitational action density

\[
\mathcal L_g=-\frac{1}{8\pi G}
\left[2\nabla\Phi\cdot\nabla\Phi_N-(g^\dagger)^2Q(z,B)\right],
\tag{10}
\]

together with the matter coupling \(-\rho_b\Phi\), gives by variation

\[
\nabla^2\Phi_N=4\pi G\rho_b
\tag{11}
\]

and

\[
\nabla^2\Phi=\nabla\cdot\left[Q_z\nabla\Phi_N\right].
\tag{12}
\]

This provides a nonrelativistic action embedding of the fixed-\(B\) response. It does not derive \(B\). If \(B\) is promoted to a dynamical field, its own kinetic/source terms, Euler--Lagrange equation, and backreaction are required. Noether conservation follows for a closed action with the relevant symmetries and all dynamical fields varied; it cannot be inferred by inserting a value of \(B\) from observed or predicted kinematics after variation.

\subsection{Algebraic rotation-curve approximation}\label{algebraic-rotation-curve-approximation}

The rotation-curve calculation uses

\[
g_{\rm eff}\simeq\Sigma g_N,\qquad
V_\Sigma\simeq V_{\rm bar}\sqrt{\Sigma}.
\tag{13}
\]

This relation is not exact for a general nonspherical mass distribution because Equation (12) contains a curl-field contribution. We compared the algebraic calculation with numerical solutions of the axisymmetric QUMOND equation for analytic disk reconstructions representative of F574-2, UGC05716, and NGC3741. Median absolute velocity differences are 5.19\%, 4.88\%, and 3.96\%, respectively; maximum local differences are 20.54\%, 18.30\%, and 7.73\%. These reconstructed models do not include the full observed gas and bulge maps, so the comparison is an approximation-error diagnostic rather than a validation of exact SPARC geometries.

\subsection{SPARC sample and statistics}\label{sparc-sample-and-statistics}

SPARC provides near-infrared photometry and resolved H I/H\(\alpha\) rotation curves for 175 disk galaxies \citep{lelli2016}. The local archive contains 171 usable rotation-curve files. Baryonic velocities are constructed with fixed mass-to-light ratios

\[
\Upsilon_{\rm disk}=0.5,\qquad
\Upsilon_{\rm bulge}=0.7,
\tag{14}
\]

and no galaxy-specific halo or response amplitude is fitted. A radial point is excluded from the primary disk analysis when the scaled bulge contribution is at least 30\% of \(V_{\rm bar}^2\). Galaxies with fewer than three remaining points are excluded. The locked sample contains 164 galaxies and 2,745 disk-dominated radial points.

For each galaxy, we calculate an unweighted velocity root-mean-square (RMS) residual for \(\Sigma\)-Gravity and for MOND. Galaxies, not radial measurements, are the independent resampling units. Uncertainty intervals use 20,000 galaxy bootstrap resamples. We use an exact binomial comparison of the fraction of galaxies for which each model has lower RMS and a paired sign-flip test of the mean RMS difference.

The frozen nuisance grid contains 81 combinations of common disk and bulge mass-to-light ratios, a global distance scale, a global inclination offset, and the \(A_0\) normalization. A separate diagnostic varies \(\sigma\) from 10 to \(50\ {\rm km\,s^{-1}}\). These are sensitivity analyses, not fitted alternatives.

\subsection{MOND comparison}\label{mond-comparison}

MOND is evaluated with the same baryonic inputs using

\[
\nu(x)=\frac{1}{1-e^{-\sqrt{x}}},
\qquad x=\frac{g_N}{a_0},
\qquad a_0=1.2\times10^{-10}\ {\rm m\,s^{-2}} .
\tag{15}
\]

The comparison is between two specified galaxy-scale phenomenological prescriptions. It is not a comparison of complete relativistic or cosmological models.

\subsection{Cluster calibration and no-refit profile check}\label{cluster-calibration-and-no-refit-profile-check}

The \citet{fox2022} catalog relates strong-lensing strength to cluster properties. The sample used here contains 42 clusters with spectroscopic redshifts and \(M_{500}>2\times10^{14}M_\odot\). The baryon approximation is

\[
M_b(<200\ {\rm kpc})=0.4\times0.15\,M_{500}.
\tag{16}
\]

The factor 0.4 is a simplified baryon-concentration prescription rather than a measured gas-plus-stellar profile and provides a reproducible baseline for the Fox calibration. The calibrated cluster amplitude is \(B_{\rm Fox}=8.446\). Writing this amplitude through a power law with \(L=600\) kpc for every cluster does not identify a path-length exponent because \(L\) does not vary within the sample. Repeated train/test splits within Fox refit the same catalog and are described only as calibration-stability checks.

The no-refit radial check uses the \citet{tian2020} CLASH radial-acceleration catalog. After removing three clusters with obvious Fox overlap, it contains 73 measurements in 17 clusters, with quoted uncertainties in \(\log g_{\rm bar}\) and \(\log g_{\rm tot}\). The Fox-calibrated \(B_{\rm Fox}\) is frozen before comparison. Residual summaries group repeated radii by cluster. Refitting this catalog is reported only as a diagnostic of catalog dependence.

\subsection{Counterrotation diagnostic}\label{counterrotation-diagnostic}

The counterrotator catalog is based on the MaNGA sample of \citet{bevacqua2022}, and the secondary dynamical quantities are drawn from MaNGA DynPop (Zhu et al., 2023). Sixty-two counterrotating systems with complete matching fields were matched to 310 controls on stellar mass, physical size, Sérsic index, axis ratio, inclination, redshift, and JAM fit quality. Balance was assessed with the absolute standardized mean difference (SMD), using 0.1 as the prespecified threshold.

The compared outcome, \(f_{\rm DM}(<R_e)\), is inferred within a Newtonian JAM/NFW model. It is not a direct observable and is retained only as a secondary diagnostic. A direct test would require common forward modeling of MaNGA velocity and dispersion maps for Newtonian, acceleration-only, and acceleration-plus-order models, with complete galaxies held out.

\subsection{Dataset roles and claim boundaries}\label{dataset-roles-and-claim-boundaries}

Table~\ref{tab:roles} identifies which results are calibration, evaluation, or sensitivity analyses. Radial measurements or spaxels within the same object are not treated as independent galaxies or clusters. The nonrelativistic action does not determine photon propagation, gravitational slip, post-Newtonian parameters, gravitational-wave propagation, or cosmological perturbations. Comparisons with lensing masses therefore require an empirical closure and are not described as derived lensing predictions.

\section{Results}\label{results}

\subsection{SPARC comparison}\label{sparc-comparison}

For 164 galaxies and 2,745 disk-dominated radial points, the mean galaxy RMS is \(16.366\ {\rm km\,s^{-1}}\) for \(\Sigma\)-Gravity and \(16.056\ {\rm km\,s^{-1}}\) for MOND. The corresponding pooled point-level root-mean-square errors are \(20.952\) and \(20.363\ {\rm km\,s^{-1}}\).

Define the paired galaxy contrast

\[
\Delta_i={\rm RMS}_{\Sigma,i}-{\rm RMS}_{{\rm MOND},i}.
\tag{17}
\]

Its galaxy-weighted mean is \(+0.309\ {\rm km\,s^{-1}}\), with bootstrap 95\% interval \([-0.040,+0.659]\ {\rm km\,s^{-1}}\). \(\Sigma\)-Gravity has lower RMS for 71 of 164 galaxies, or 43.3\%, with bootstrap interval 36.0\%--50.6\%. The exact two-sided binomial \(p\)-value is 0.101, and the paired sign-flip \(p\)-value is 0.088. The models therefore show comparable aggregate performance; neither the paired interval nor the object-level fraction establishes superiority.

The conclusion is sensitive to common nuisance choices. Across the 81 frozen combinations, the mean paired contrast ranges from \(-3.286\) to \(+2.779\ {\rm km\,s^{-1}}\) and retains the central sign in 64.2\% of cases (Figure~\ref{fig:1}). This dependence is a reason not to overinterpret the small difference between the two prescriptions.

\subsection{Endogenous-factor ablation}\label{endogenous-factor-ablation}

Replacing Equation (6) with the acceleration-only form \(B=A_0\) gives a mean galaxy RMS of \(16.882\ {\rm km\,s^{-1}}\). The locked endogenous factor therefore improves the mean RMS by \(0.517\ {\rm km\,s^{-1}}\), with bootstrap 95\% interval \(0.242\)--\(0.795\ {\rm km\,s^{-1}}\). This establishes that the factor changes the predictor; it does not establish an independent coherence effect because \(F\) is calculated from \(V_\Sigma\). The correct interpretation is an internal ablation of the phenomenological equation.

\subsection{Algebraic approximation error}\label{algebraic-approximation-error}

The numerical QUMOND reconstructions show that the algebraic relation in Equation (13) has a radius-dependent error (Figure~\ref{fig:3}). The median absolute discrepancy is approximately 4\%--5\% in the three representative disks, while the largest inner-radius deviation reaches 20.5\%. This uncertainty is comparable to or larger than the aggregate RMS separation between \(\Sigma\)-Gravity and MOND for some galaxies. It should be incorporated into future object-level likelihoods rather than treated as an exact identity.

\subsection{Fox cluster calibration}\label{fox-cluster-calibration}

Under Equation (16), calibrating the common cluster response on the 42 Fox systems gives a median predicted-to-observed mass ratio of 0.987 and a scatter of 0.132 dex (Figure~\ref{fig:2}, left). Because the amplitude was selected using these targets, the ratio is an in-sample calibration result. Repeated Fox holdouts measure stability under resampling of the same catalog; they do not constitute independent validation.

The baryon prescription is materially influential. Varying the 0.4 concentration factor by 25\% changes the predicted mass ratio by approximately 30\%. The calibrated amplitude cannot therefore be separated from the assumed baryonic aperture mass.

\subsection{No-refit CLASH profile check}\label{no-refit-clash-profile-check}

With \(B_{\rm Fox}=8.446\) held fixed, the 73 measurements in 17 non-overlapping CLASH clusters give a median predicted-to-observed ratio of 1.318 and an RMS residual of 0.188 dex. Median ratios are 1.18 at 100 kpc, 1.34 at 200 kpc, 1.66 at 400 kpc, and 1.98 at 600 kpc (Figure~\ref{fig:2}, right). Thus the Fox-calibrated amplitude increasingly overpredicts the profile at larger radius under the present baryonic and lensing assumptions.

This no-refit result does not identify a replacement amplitude or a revised cluster formula. It shows that the Fox-calibrated amplitude does not transfer without radial bias under the present assumptions.

\subsection{Counterrotation}\label{counterrotation}

Matching reduces the maximum absolute SMD from 0.684 to 0.071, below the 0.1 balance threshold (Figure~\ref{fig:4}). The matched mean difference is

\[
\Delta f_{\rm DM}
=f_{{\rm DM},{\rm counter}}-f_{{\rm DM},{\rm control}}
=-0.0069 ,
\tag{18}
\]

with bootstrap 95\% interval \([-0.0567,+0.0466]\). The result is consistent with zero and does not reproduce the large unmatched association. Because \(f_{\rm DM}\) is JAM/NFW-derived, the matched result is not a direct test of \(\Sigma\)-Gravity. Counterrotation remains a proposed future discriminant rather than a confirmed prediction.

\section{Discussion}\label{discussion}

\subsection{A bounded alternative}\label{a-bounded-alternative}

\(\Sigma\)-Gravity is an alternative empirical response law for selected missing-mass observations. At galaxy scale, it replaces a fitted halo or a MOND interpolation function with a specified baryon-dependent prescription. The locked SPARC result shows that this low-parameter prescription remains close to MOND in aggregate without fitting an individual response amplitude to each galaxy. That is an interesting empirical result, but it is not evidence that dark matter is absent or that the model supersedes MOND.

At cluster scale, the current implementation is weaker. The Fox result shows that a common response can be calibrated under a simple baryon prescription. The no-refit CLASH profiles show that this normalization does not transfer without radial bias under the present assumptions. We therefore retain the Fox value as a calibration baseline and do not introduce a replacement cluster formula.

\subsection{Coherence and source organization as hypotheses}\label{coherence-and-source-organization-as-hypotheses}

The empirical response does not identify its physical mechanism. One possibility is that the effective response depends on the organization or dynamical state of the baryonic source. Kinematic coherence or another correlated property could contribute to a dynamical order parameter.

For independently measured phase-space data, one operational estimator would be

\[
\mathcal C_{\rm kin}=
\frac{|\bar{\mathbf v}_{\rm stream}|^2}
{|\bar{\mathbf v}_{\rm stream}|^2+
{\rm tr}\,\boldsymbol{\sigma}^2}
\tag{19}
\]

in the system barycentric frame. This quantity would be small for a dispersion-supported cluster, whereas the Fox calibration uses a single effective response rather than an independently measured \(\mathcal C_{\rm kin}\). Therefore one measured coherence definition does not currently explain both disks and clusters. The hypothesis should advance only if independently measured phase-space order improves held-out predictions relative to an acceleration-only model, with a stable sign across mass and morphology.

\subsection{Path length and source concentration}\label{path-length-and-source-concentration}

A path-length power law can compactly relate galaxy and cluster amplitudes, but no unique functional for a general three-dimensional baryonic distribution is defined here. Because every Fox cluster received \(L=600\) kpc, the sample cannot identify the exponent \(n\) independently. Path length is therefore retained as a possible motivation for future work, not as a derived part of the canonical model.

\subsection{Relation to Refracted Gravity}\label{relation-to-refracted-gravity}

A related but distinct future direction is Refracted Gravity, in which a density-dependent gravitational permittivity modifies the Poisson equation and can redirect field lines in nonspherical sources \citep{cesare2020,cesare2022,sanna2023}. This literature shows how density and source geometry might enter a more developed theory, but it does not derive or validate the empirical \(\Sigma\) response. Any connection would require a separate theoretical and observational study.

\subsection{Lensing, Solar-System, and relativistic limitations}\label{lensing-solar-system-and-relativistic-limitations}

In a general weak-field metric,

\[
ds^2=-(1+2\Psi)c^2dt^2+(1-2\Phi)d\mathbf x^2,
\tag{20}
\]

nonrelativistic dynamics primarily constrains \(\Psi\), while lensing depends on \(\Phi+\Psi\). The gravitational slip \(\eta=\Phi/\Psi\) is not determined by Equations (10)--(12). Equality between a phenomenological dynamical mass and a lensing mass is therefore a closure assumption, not a derived result.

At Solar-System accelerations, \(h(g_N)\) is small. This is a useful anomalous-acceleration sanity check, but it is not a calculation of the parameterized post-Newtonian parameter \(\gamma\), and the assumption that a compact source has no additional response has not been derived. Equivalence-principle, post-Newtonian, and gravitational-wave constraints remain open pending a relativistic completion.

\subsection{Cosmological scope}\label{cosmological-scope}

No cosmological extension is developed here, and the present response law should not be interpreted as an alternative cosmology. Any such extension would need to address the cosmic microwave background spectrum, primordial abundances, perturbation growth, nonlinear structure formation, and cluster-abundance evolution that form central parts of the \(\Lambda\)CDM evidence base.

\subsection{Decisive next tests}\label{decisive-next-tests}

The highest-value tests add independent observables rather than formula flexibility. First, MaNGA velocity and dispersion maps should be forward-modeled for counterrotators and balanced controls under common baryonic and nuisance assumptions, with complete galaxies held out. Second, a disjoint cluster sample should combine radial X-ray gas, brightest-cluster-galaxy, satellite, and intracluster-light profiles with lensing covariance fixed before any amplitude is inferred. Third, exact axisymmetric or three-dimensional field solutions should replace the algebraic disk approximation in the model likelihood.

An acceleration-plus-coherence model should advance only if it improves galaxy-grouped held-out performance over acceleration-only predictions by a prespecified margin and remains stable across morphology and mass. A cluster law should be called predictive only if it transfers without refitting after the baryon and lensing likelihoods are frozen.

\section{Conclusions}\label{conclusions}

\(\Sigma\)-Gravity provides a compact acceleration-based phenomenology for selected missing-mass observations. On a locked sample of 164 SPARC galaxies, its aggregate rotation-curve performance is statistically comparable to the MOND prescription tested with the same fixed baryonic assumptions and without fitting an individual halo or response amplitude to each galaxy. Within the \(\Sigma\)-Gravity predictor, the bounded rotational-support-motivated factor improves the mean RMS over the acceleration-only form by \(0.517\ {\rm km\,s^{-1}}\), with a 95\% interval of \([0.242,0.795]\ {\rm km\,s^{-1}}\). This is a statistically supported internal contribution, although its endogenous construction does not establish independently measured coherence.

The Fox cluster result is an in-sample calibration conditional on a simplified baryon model and an assumed relation to lensing mass. A no-refit CLASH profile check reveals an increasing radial overprediction. The calibrated \(B=8.45\) is therefore not established as a universal cluster value, and a path-length interpretation remains a hypothesis.

An action-based nonrelativistic QUMOND embedding exists for independently specified \(B\), but the physical origin and dynamics of that amplitude remain unknown. Coherence or another property of baryonic organization is a possible explanation worth testing; it is not established by the current data. The matched counterrotation result is consistent with zero and is reported as a negative secondary diagnostic.

Together, these results establish a fixed low-parameter galaxy response with a measurable incremental contribution, a QUMOND action embedding for independently specified \(B\), and an external cluster test that rejects unbiased transfer of the fixed Fox-calibrated amplitude under the present assumptions. \(\Sigma\)-Gravity therefore merits continued investigation as a galaxy-scale phenomenological alternative. Direct phase-space measurements and independently measured cluster baryon profiles are required to determine whether the empirical response reflects new physics or a compact description of correlated astrophysical structure; no relativistic or cosmological extension is developed here.

\section*{Data Availability Statement}

SPARC data are publicly available at \href{http://astroweb.cwru.edu/SPARC/}{astroweb.cwru.edu/SPARC}. The CLASH radial-acceleration catalog is available through VizieR as J/ApJ/896/70. The Fox cluster table used for calibration, frozen residuals, split definitions, matched samples, parameter diagnostics, and figure-generation code are provided in the accompanying repository and Supplementary Material.

\section*{Author Contributions}

LS conceived the study, developed the model, assembled the data and code, performed the analyses, interpreted the results, and wrote and revised the manuscript.

\section*{Funding}

The author declares that no financial support was received for the research, authorship, and/or publication of this article.

\section*{Conflict of Interest Statement}

The author declares that the research was conducted in the absence of any commercial or financial relationships that could be construed as a potential conflict of interest.

\section*{Acknowledgments}

The author thanks the reviewers and editor for comments that materially improved the scope, statistical presentation, and claim boundaries of the manuscript.

\section*{Supplemental Data}

The Supplementary Material contains the full statistical procedures, identifiability audit, frozen dataset roles, software-regression commands, and machine-readable output manifest.

\bibliographystyle{Frontiers-Harvard}
\bibliography{SigmaGravity_Resubmission_2026-07-25}

\section*{Tables}

\begin{table}[h!]
\centering
\caption{Role of each dataset and analysis.}\label{tab:roles}
\scriptsize
\begin{tabularx}{\textwidth}{p{0.20\textwidth}p{0.25\textwidth}p{0.20\textwidth}X}
\toprule
Dataset or analysis & Role & Independent unit & Determines a model parameter? \\
\midrule
SPARC disk sample & Locked cross-domain empirical evaluation & Galaxy & No per-galaxy parameter \\
Fox clusters & Calibration & Cluster & Yes, one effective cluster amplitude \\
Repeated Fox splits & Calibration-stability diagnostic & Held-out cluster within Fox & Refits on each training subset \\
Non-overlapping Tian/CLASH profiles & No-refit external profile check & Cluster, with radii grouped & No \\
Matched MaNGA/JAM catalog & Secondary counterrotation diagnostic & Counterrotator/control set & No \\
Numerical QUMOND disks & Algebraic-approximation diagnostic & Reconstructed galaxy model & No \\
\bottomrule
\end{tabularx}
\end{table}

\begin{table}[h!]
\centering
\caption{Parameter and assumption accounting.}\label{tab:parameters}
\scriptsize
\begin{tabularx}{\textwidth}{p{0.18\textwidth}p{0.22\textwidth}p{0.18\textwidth}X}
\toprule
Quantity & Role & Status & Principal limitation \\
\midrule
\(g^\dagger\) & Acceleration scale in \(h(g_N)\) & Fixed model choice & BTFR constrains \(B^2g^\dagger\), not each factor \\
\(A_0=e^{1/(2\pi)}\) & Galaxy normalization before \(F(V_\Sigma)\) & Fixed model choice & Not uniquely derived from the data \\
\(\sigma=20\ {\rm km\,s^{-1}}\) & Regulator in \(F(V_\Sigma)\) & Fixed & Endogenous; sensitivity tested from 10--50 km s\(^{-1}\) \\
\(\Upsilon_{\rm disk},\Upsilon_{\rm bulge}\) & Stellar baryonic contributions & Fixed astrophysical assumptions & Change relative \(\Sigma\)/MOND performance \\
\(B_{\rm Fox}=8.446\) & Cluster response & Calibrated & Does not transfer as a universal radial amplitude \\
\(0.4\times0.15M_{500}\) & Fox baryon mass inside 200 kpc & Approximation & Material amplitude sensitivity \\
Lensing closure & Relates response to lensing target & Assumed & Gravitational slip is undetermined \\
\(R_d,L_0,n\) & Scale-length/path hypothesis & Not used in locked galaxy result & No unique operational 3D functional \\
\bottomrule
\end{tabularx}
\end{table}

\section*{Figure captions}

\begin{figure}[h!]
\begin{center}
\includegraphics[width=0.98\textwidth]{figures/figure_1_sparc_paired.pdf}
\end{center}
\caption{Locked SPARC comparison and nuisance sensitivity. Left: per-galaxy velocity RMS for \(\Sigma\)-Gravity and the tested MOND prescription, with SPARC quality classes shown separately. Center: distribution of the paired contrast \({\rm RMS}_\Sigma-{\rm RMS}_{\rm MOND}\); the mean is \(+0.309\ {\rm km\,s^{-1}}\) and its 95\% galaxy-bootstrap interval includes zero. Right: mean contrast for the 81 frozen nuisance combinations. Negative values favor \(\Sigma\)-Gravity and positive values favor MOND. The grid spans both signs, so the models are described as having comparable aggregate performance.}\label{fig:1}
\end{figure}

\begin{figure}[h!]
\begin{center}
\includegraphics[width=0.98\textwidth]{figures/figure_2_cluster_roles.pdf}
\end{center}
\caption{Cluster calibration and no-refit radial evaluation. Left: predicted and observed 200-kpc aperture masses for the 42 Fox clusters under the simplified baryon proxy and calibrated \(B_{\rm Fox}\). This panel is an in-sample calibration, not validation. Right: predicted-to-observed acceleration ratios for 73 Tian/CLASH radial measurements in 17 clusters after obvious Fox overlaps are removed, with \(B_{\rm Fox}\) frozen. Large markers show radius-bin medians. The systematic increase with radius shows that the fixed Fox-calibrated amplitude does not transfer without bias under the present baryonic and lensing assumptions.}\label{fig:2}
\end{figure}

\begin{figure}[h!]
\begin{center}
\includegraphics[width=0.98\textwidth]{figures/figure_3_qumond_approximation.pdf}
\end{center}
\caption{Algebraic approximation error in representative axisymmetric disk reconstructions. Fractional velocity difference between the algebraic relation in Equation (13) and the numerical QUMOND solution for analytic reconstructions representative of F574-2, NGC3741, and UGC05716. The comparison quantifies a geometry-dependent approximation error; it is not a fit to the full observed gas and bulge maps.}\label{fig:3}
\end{figure}

\begin{figure}[h!]
\begin{center}
\includegraphics[width=0.98\textwidth]{figures/figure_4_counterrotation_matched.pdf}
\end{center}
\caption{Matched counterrotation diagnostic. Left: absolute standardized mean differences before and after matching counterrotators to controls. All post-match values are below the prespecified 0.1 threshold. Right: matched difference in the JAM/NFW-derived \(f_{\rm DM}(<R_e)\), with galaxy-bootstrap 95\% interval. The interval includes zero. Because the outcome is model derived, this is a secondary catalog diagnostic rather than a direct test of \(\Sigma\)-Gravity.}\label{fig:4}
\end{figure}

\end{document}
